\scanpage{133}

\mspage{469}

d'eux et à lui-même, et aux deux autres: conflit insidieux,au visage impassible entre ma mère et ma soeur, et confli\unsure{t} aux violents éclats entre mon père et ma mère d'un côté, ma soeur de l'autre, qui chacune pour son propre compte (et sans que personne du vivant de mes parents ait jamais fait mine de s'en apercevoir... le faisait marcher à sa foçon. La chose mystérieuse, extraordinaire, c'est qu'entouré ainsi par le conflit en ces années les plus sensibles, les plus cruciales de la vie, celui-ci soit resté \mbox{e x t é r i e u r} à moi, qu'il n'ait pas vraiment "mordu" sur mon être en ces années-là et ne s'y soit installé à demeure.

La division dans mon être, qui a marqué ma vie tout autant que celle de \unsure{to|tout} autre, ne s'est pas installée en moi en ces années, mais en les deux ou trois \unsure{an|ans} qui ont suivi, de ma sixième à ma huitième année environ. A un certain moment. (que j'ai cru pouvoir situer à quelques mois près, et qui se placerait dans ma huitième année) il y a eu un certain \mbox{b a s c u l e m e n t}, au bout de plus de deux ans de séparation avec mes parents (qui ne se souciaient guère de me donner signe de vie) et de ma soeur. C'était avant tout une \mbox{r u p t u r e}\quad\mbox{a v e c}\quad\mbox{m o n}\quad\mbox{e n f a n c e}, "enterrée" à partir de ce moment par d'efficaces mécanismes d'oubli (qui sont restés en place, à peu de choses près, jusqu'à aujourd'hui même). A un certain niveau profond (pas le plus profond pourtant...) mes parents ont alors été déclarés par moi comme des "étrangers", tout comme mon enfance était désormais déclarée "étrangère". J'ai \mbox{a b\unsure{.}d i q u é}, en un sens: pour être accepté dans le monde qui m'entourait désormais, j'ai décidé d'être comme "eux", comme les adultes qui y font la loi - d'acquérir et de développer les armes qui y forcent le respect, de me battre à armes égales dans un monde où seule, une certaine espèce de "force" est acceptée et prisée...

C'était d'ailleurs cette force-là aussi qui avait la préférence de mes parents, lesquels avaient entouré mes premières années. Et là je reviens à cette "chose mystérieuse" (dont je viens de m'éloigner, en suivant le fil d'une autre association suscitée par cette chose), \mbox{l' a b s e n c e}\quad\mbox{d e}\quad\mbox{d i v i s i o n}\quad\mbox{e n}\quad\mbox{m o i}, en ces premières années de ma vie.

Peut-être le mystère n'est plus pour moi en cette absence, mais plutôt en ceci: que mes parents, mon père comme ma mère, m'aient chacun alors \mbox{\unsure{a c c e p t|a c c e p t é}}\quad\mbox{d a n s}\quad\mbox{m a}\quad\mbox{t o t a l i t é}, et totalement: dans ce qui en moi est "viril" est "homme", et dans ce qui est "femme". Ou pour le dire autrement: que mes parents, déchirés l'un et l'autre par le conflit, reniant chacun une partie essentielle de leur être - incapable chacun d'une ouverture aimante à lui-même et à l'autre, comme d'une ouverture aimante à ma soeur... que néanmoins ils aient trouvé une telle ouverture, une acceptation sans réserve, vis à vis de moi

% archive sheet number: 132
